\documentclass[11pt, a4paper]{article}

% --- ESSENTIAL PACKAGES ---
\usepackage[utf8]{inputenc}
\usepackage[T1]{fontenc}
\usepackage[english]{babel}

% --- MATHEMATICS ---
\usepackage{amsmath}
\usepackage{amssymb}
\usepackage{mathtools} % For \coloneqq, etc.

% --- LAYOUT ---
\usepackage{geometry}
\geometry{a4paper, margin=1in}

% --- HYPERLINKS ---
\usepackage{hyperref}
\hypersetup{
    colorlinks=true,
    linkcolor=blue,
    filecolor=magenta,      
    urlcolor=cyan,
    pdftitle={Complete Summary of the Technical Discussion},
    pdfauthor={Janis STG3A MMNN},
}

% --- CUSTOM COMMANDS ---
\newcommand{\E}{\mathbb{E}}
\newcommand{\dd}{\,d} % upright 'd' for integrals
\DeclareMathOperator{\ReLU}{ReLU}
\DeclareMathOperator{\Var}{Var}
\newcommand{\indic}{\mathbb{1}}

\title{Complete Summary of the Technical Discussion}
\author{Janis STG3A MMNN}
\date{\today}

\begin{document}

\maketitle

\begin{abstract}
This document recounts all the calculations, derivations, and concepts covered during our technical discussion. It covers the calculation of the expectation of products of ReLU functions applied to bivariate Gaussian variables, the application of Price's theorem, the analysis of special cases, and the study of the underlying probability laws such as the Fisher and Kibble distributions, in relation to the curse of dimensionality.
\end{abstract}

\section{Calculating the Expectation $\E[\ReLU(X_1)\ReLU(X_2)]$}
The initial problem is to calculate the expectation $\E[\ReLU(X_1)\ReLU(X_2)]$ for a pair of random variables $(X_1, X_2)$ following a bivariate normal distribution with parameters $(\mu_1, \sigma_1, \mu_2, \sigma_2, \rho)$.

\subsection{Decomposition}
The approach is to standardize the variables ($Z_i = (X_i - \mu_i)/\sigma_i$) and decompose the integral into four terms:
\[
\E[\dots] = \sigma_1\sigma_2 I_{22} + \mu_1\sigma_2 I_{12} + \mu_2\sigma_1 I_{21} + \mu_1\mu_2 I_{11}
\]
where the integrals $I_{ij}$ are defined over the domain $z_1 > a$ and $z_2 > b$, with $a = -\mu_1/\sigma_1$ and $b = -\mu_2/\sigma_2$.
\begin{itemize}
    \item $I_{11} = \int_a^\infty \int_b^\infty \phi_2(z_1, z_2; \rho) \dd z_2 \dd z_1$
    \item $I_{21} = \int_a^\infty \int_b^\infty z_1 \phi_2(z_1, z_2; \rho) \dd z_2 \dd z_1$
    \item $I_{12} = \int_a^\infty \int_b^\infty z_2 \phi_2(z_1, z_2; \rho) \dd z_2 \dd z_1$
    \item $I_{22} = \int_a^\infty \int_b^\infty z_1 z_2 \phi_2(z_1, z_2; \rho) \dd z_2 \dd z_1$
\end{itemize}

\subsection{Calculation of the Components}

\subsubsection*{Calculation of $I_{11}$ (Probability)}
$I_{11}$ is the probability $P(Z_1 > a, Z_2 > b)$. It is calculated via the cumulative distribution function (CDF) of the standard bivariate normal distribution, $\Phi_2$.
\[
I_{11} = P(Z_1 > a, Z_2 > b) = P(Z_1 < -a, Z_2 < -b) = \Phi_2(-a, -b; \rho) = \Phi_2\left(\frac{\mu_1}{\sigma_1}, \frac{\mu_2}{\sigma_2}; \rho\right)
\]

\subsubsection*{Calculation of $I_{22}$ (Cross-Moment)}
The calculation of $I_{22}$ is the most complex part. The final result is:
\[
I_{22} = \rho I_{11} + \phi_2(a, b; \rho)
\]
The proof of this identity was a central point of the discussion. The most rigorous method consists of showing the equality for $\rho=0$, and then the equality of the derivatives with respect to $\rho$ for both sides of the equation.

\subsubsection*{Final Formula}
By assembling all the terms, we obtain the final algebraic formula:
\[
\E[\dots] = (\mu_1\mu_2 + \rho\sigma_1\sigma_2) I_{11} + \mu_1\sigma_2 I_{12} + \mu_2\sigma_1 I_{21} + \sigma_1\sigma_2 \phi_2(a, b; \rho)
\]

\section{Price's Theorem and Derivations}
Price's theorem was a central tool in our analysis.

\subsection{Differential and Integral Form}
The theorem relates the derivative of the expectation of a function $g$ with respect to the covariance $\sigma_{12}$ to the expectation of the second derivative of $g$. For $g = \ReLU(X_1)\ReLU(X_2)$, we have $\frac{\partial^2 g}{\partial x_1 \partial x_2} = \indic(X_1>0, X_2>0)$.
\[
\frac{\partial \E[\ReLU(X_1)\ReLU(X_2)]}{\partial \sigma_{12}} = \E[\indic(X_1>0, X_2>0)] = I_{11}
\]
This leads to the integral formulation:
\[
\E[\ReLU(X_1)\ReLU(X_2)] = \int I_{11}(\sigma_{12}) \dd\sigma_{12} + \text{Constant}
\]

\subsection{Calculating the Constant (Independent Case $\rho=0$)}
The constant is the value of the expectation when $\rho=0$, i.e., when $X_1$ and $X_2$ are independent.
\[
C = \E[\ReLU(X_1)] \cdot \E[\ReLU(X_2)]
\]
The expectation of a single ReLU is:
\[
\E[\ReLU(X)] = \sigma\phi_1\left(\frac{\mu}{\sigma}\right) + \mu\Phi_1\left(\frac{\mu}{\sigma}\right)
\]
The constant is therefore the product of two such terms:
\begin{multline*}
C = \left[ \sigma_1\phi_1\left(\frac{\mu_1}{\sigma_1}\right) + \mu_1\Phi_1\left(\frac{\mu_1}{\sigma_1}\right) \right] \times \\ \left[ \sigma_2\phi_1\left(\frac{\mu_2}{\sigma_2}\right) + \mu_2\Phi_1\left(\frac{\mu_2}{\sigma_2}\right) \right]
\end{multline*}

\section{Unsuccessful Attempts to Calculate $I_{22}$}
Several approaches to calculate $I_{22}$ proved to be dead ends.

\subsection{Via the Identity on the Derivative of $\phi_2$}
The identity $\int_a^{\infty} z_1 \phi_2 \dd z_1 = \rho z_2 \int_a^{\infty} \phi_2 \dd z_1 + (1-\rho^2)\phi_2(a, z_2)$ has been rigorously proven. An attempt to use it to calculate $I_{22}$ leads to the expression:
\[
I_{22} = \rho \int_a^\infty z_1^2 \left(\int_b^\infty \phi_2 \dd z_2\right) \dd z_1 + (1-\rho^2) \int_a^\infty z_1 \phi_2(z_1, b; \rho) \dd z_1
\]
This approach fails because the first term contains a second-order moment ($z_1^2$), which is more complex than the initial problem, leading to increasing complexity.

\subsection{Via Decomposition into 1D PDF}
Another attempt was to use the decomposition $\phi_2(z_1, z_2) = \phi_1(z_1) f(z_2|z_1)$. Applied to the terms above, this method leads to integrals with no simple analytical solution.

\section{Special Case: Zero Means ($\mu_1=\mu_2=0$)}
In this case, the "from scratch" derivation leads to an elegant formula.
\[
\E[\ReLU(X_1)\ReLU(X_2)] = \sigma_1\sigma_2 \int_0^\infty \int_0^\infty z_1 z_2 \phi_2(z_1, z_2; \rho) \dd z_1 \dd z_2
\]
The problem reduces to calculating $I_{22}(0,0; \rho) = \E[\ReLU(Z_1)\ReLU(Z_2)]$. The cleanest derivation uses Price's theorem:
\begin{enumerate}
    \item $\frac{\partial \E[g]}{\partial \rho} = \E\left[\frac{\partial^2 g}{\partial z_1 \partial z_2}\right] = \E[\indic(z_1>0, z_2>0)] = P(Z_1>0, Z_2>0) = \frac{1}{2\pi}\arccos(-\rho)$.
    \item The integration constant is $\E[g]_{\rho=0} = \E[\ReLU(Z_1)]\E[\ReLU(Z_2)] = \left(\frac{1}{\sqrt{2\pi}}\right)^2 = \frac{1}{2\pi}$.
    \item By integrating $\frac{\partial \E[g]}{\partial \rho}$ from $0$ to $\rho$:
    \begin{align*}
        \E[g](\rho) - \E[g](0) &= \int_0^\rho \frac{1}{2\pi}\arccos(-t) \dd t \\
        &= \frac{1}{2\pi} \left[ t\arccos(-t) + \sqrt{1-t^2} \right]_0^\rho \\
        &= \frac{1}{2\pi} (\rho\arccos(-\rho) + \sqrt{1-\rho^2} - 1)
    \end{align*}
    \item By adding the constant $\E[g](0) = 1/(2\pi)$, we get:
    \[
    \E[\ReLU(Z_1)\ReLU(Z_2)] = \frac{1}{2\pi}(\sqrt{1-\rho^2} + \rho\arccos(-\rho))
    \]
\end{enumerate}
The final formula for variances $\sigma_1^2, \sigma_2^2$ is:
\[
\E[\ReLU(X_1)\ReLU(X_2)] = \frac{\sigma_1\sigma_2}{2\pi}(\sqrt{1-\rho^2} + \rho\arccos(-\rho))
\]

\section{Fisher, Kibble distributions and the "Curse of Dimensionality"}
The discussion then shifted to the underlying probability laws of the random components of the problem.

\subsection{The Fisher Distribution for the Correlation Coefficient $r$}
\paragraph{Distribution} The exact law of $r$ (cosine of the angle between two Gaussian vectors of dimension $n$ with correlation $\rho$) is complex.
\begin{itemize}
    \item \textbf{Reference:} R. A. Fisher (1915), "Frequency Distribution...", \textit{Biometrika}.
    \item \textbf{Modern Form:}
        \[
        p(r|\rho, n) = \frac{(n-2)\Gamma(n-1)(1-\rho^2)^{\frac{n-1}{2}}(1-r^2)^{\frac{n-4}{2}}}{\sqrt{2\pi}\Gamma(n-\frac{1}{2})(1-\rho r)^{n-\frac{3}{2}}} \times {}_2F_1\left(\frac{1}{2}, \frac{1}{2}; n-\frac{1}{2}; \frac{1+\rho r}{2}\right)
        \]
\end{itemize}

\paragraph{Approximation (Fisher's z-transformation)} For large $n$, $z = \text{arctanh}(r)$ follows a $\mathcal{N}\left(\text{arctanh}(\rho) + \frac{\rho}{2(n-1)}, \frac{1}{n-3}\right)$ distribution.

\paragraph{Curse of Dimensionality} The variance of $r$ decreases as $O(1/n)$, not exponentially. If $\rho \sim O(1/\sqrt{n})$, the correlation signal is drowned out by geometric noise.

\subsection{The Kibble Distribution for the Squared Norms $(U,V)$}
\paragraph{Distribution} The joint law of $U=\|X\|^2 \sim \chi^2(n)$ and $V=\|Y\|^2 \sim \chi^2(n)$ (correlated) is the bivariate Kibble-Gamma distribution.
\begin{itemize}
    \item \textbf{Reference:} W. F. Kibble (1941), "A Two-Variate Gamma Type Distribution", \textit{Sankhyā}.
    \item \textbf{Formula:} 
    \[
    f(u, v) = \frac{(uv)^{\frac{n/2-1}{2}} e^{-\frac{u+v}{2(1-\rho^2)}}}{\Gamma(n/2) (2(1-\rho^2))^{n/2+1} \rho^{\frac{n/2-1}{2}}} \cdot I_{n/2-1}\left(\frac{\rho\sqrt{uv}}{1-\rho^2}\right)
    \]
\end{itemize}

\paragraph{Curse of Dimensionality}
\begin{itemize}
    \item $\E[U] = n$, $\Var(U) = 2n$. The mean and variance increase.
    \item The coefficient of variation $CV = \sqrt{2n}/n = \sqrt{2}/\sqrt{n}$ tends to 0. The distribution becomes tighter relative to its mean.
    \item The normalized variable $U/n$ has a mean of $1$ and variance of $2/n$. It converges to the constant 1 (Law of Large Numbers) following a $\mathcal{N}(1, 2/n)$ approximation (Central Limit Theorem).
\end{itemize}

\subsection{Product of Kibble Variables $W=UV$}
\begin{itemize}
    \item \textbf{Law:} No simple name. Expressed via infinite series or special functions (Meijer-G, Whittaker).
    \item \textbf{Mean:} $\E[W] = n^2 + 2n\rho^2$.
    \item \textbf{Variance:} $\Var(W) = (4n^3 + 4n^2) + (4n^3 + 32n^2 + 32n)\rho^2 + (4n^2 + 16n)\rho^4$.
    \item \textbf{Variance of the normalized product $W/n^2$:}
    \[
    \Var\left(\frac{W}{n^2}\right) = \left(\frac{4}{n} + \frac{4}{n^2}\right) + \left(\frac{4}{n} + \frac{32}{n^2} + \frac{32}{n^3}\right)\rho^2 + \left(\frac{4}{n^2} + \frac{16}{n^3}\right)\rho^4 \approx \frac{4(1+\rho^2)}{n}
    \]
\end{itemize}

\subsection{Independence}
A key result is that the variable $r$ (information on relative orientation) is \textbf{independent} of the pair $(U,V)$ (information on magnitudes).

\subsection{Link with the Gudermannian Function}
The Fisher transformation and the angle $\theta=\arccos(r)$ are related via the Gudermannian function, but this is a mathematical curiosity rather than a simplifying tool.
\[
z = \text{arctanh}(r) = \text{gd}^{-1}\left(\frac{\pi}{2} - \theta\right)
\]

\section{Detailed Derivations from Discussion}
This section provides a more granular, step-by-step walkthrough of the calculations presented in the main body of the paper, as captured during the technical discussion.

\subsection{Step-by-Step Calculation of the Expectation}

\subsubsection{Problem Definition}
We want to compute the expectation of the product of two random variables, $X_1$ and $X_2$, truncated at zero by the ReLU function:
\[
\E[\ReLU(X_1)\ReLU(X_2)] = \E[\max(0, X_1)\max(0, X_2)]
\]
The random variables $(X_1, X_2)$ follow a bivariate normal distribution with means $(\mu_1, \mu_2)$, variances $(\sigma_1^2, \sigma_2^2)$, and correlation coefficient $\rho$. The expectation is calculated by integrating over the domain where $X_1 > 0$ and $X_2 > 0$:
\[
\E[\max(0,X_1)\max(0,X_2)] = \int_0^\infty \int_0^\infty x_1 x_2 \cdot f_{X_1,X_2}(x_1,x_2) \dd x_1 \dd x_2
\]
where $f_{X_1,X_2}$ is the PDF of the bivariate normal distribution.

\subsubsection{Change of Variables}
To simplify the calculation, we standardize the variables. Let:
\begin{align*}
    Z_1 &= \frac{X_1 - \mu_1}{\sigma_1} \implies X_1 = \sigma_1 Z_1 + \mu_1 \\
    Z_2 &= \frac{X_2 - \mu_2}{\sigma_2} \implies X_2 = \sigma_2 Z_2 + \mu_2
\end{align*}
The variables $(Z_1, Z_2)$ follow a standard bivariate normal distribution with correlation $\rho$, whose PDF is $\phi_2(z_1, z_2; \rho)$. By substituting into the integral, the integration bounds change from $x_i > 0$ to $z_i > -\mu_i/\sigma_i$. Let $a = -\mu_1/\sigma_1$ and $b = -\mu_2/\sigma_2$.

\subsubsection{Decomposition of the Integral}
By expanding the integrand and using the linearity of the integral, we get four terms:
\begin{align*}
\E[\dots] &= \int_a^\infty \int_b^\infty (\sigma_1 z_1 + \mu_1)(\sigma_2 z_2 + \mu_2) \cdot \phi_2(z_1, z_2; \rho) \dd z_1 \dd z_2 \\
&= \sigma_1\sigma_2 \underbrace{\int_a^\infty \int_b^\infty z_1 z_2 \phi_2 \dd z_1 \dd z_2}_{I_{22}} + \mu_1\sigma_2 \underbrace{\int_a^\infty \int_b^\infty z_2 \phi_2 \dd z_1 \dd z_2}_{I_{12}} \\
&\quad + \mu_2\sigma_1 \underbrace{\int_a^\infty \int_b^\infty z_1 \phi_2 \dd z_1 \dd z_2}_{I_{21}} + \mu_1\mu_2 \underbrace{\int_a^\infty \int_b^\infty \phi_2 \dd z_1 \dd z_2}_{I_{11}}
\end{align*}

\subsubsection{Calculation of the Truncated Moments}
\paragraph{$I_{11}$} This term is the probability of the variables being in the integration domain, which is the definition of the bivariate normal CDF:
\[
I_{11} = \Phi_2(-a, -b; \rho) = \Phi_2\left(\frac{\mu_1}{\sigma_1}, \frac{\mu_2}{\sigma_2}; \rho\right)
\]

\paragraph{$I_{12}$ and $I_{21}$} These are the first-order moments over the truncated domain. Their calculation involves complex integration, and the formulas provided during the discussion were:
\begin{align*}
    I_{12} &= \phi_1(a) \Phi_1(b - \rho a) \\
    I_{21} &= \phi_1(b) \Phi_1(a - \rho b)
\end{align*}
Note: These are simplified forms. The full expression for these moments is more complex. For instance, a more complete formula for the truncated expectation is $E[Z_2 | Z_1>a, Z_2>b] = \frac{1}{I_{11}} \left( \phi_1(a)\Phi_1(\frac{\rho a - b}{\sqrt{1-\rho^2}}) + \rho\phi_1(b)\Phi_1(\frac{\rho b - a}{\sqrt{1-\rho^2}}) \right)$.

\paragraph{$I_{22}$} The calculation of the second-order cross-moment is the most subtle part. It relies on a remarkable identity of the bivariate normal distribution, related to Price's Theorem, which connects the derivative of the CDF ($I_{11}$) with respect to the correlation $\rho$ to the PDF $\phi_2$:
\[
\frac{\partial I_{11}}{\partial \rho} = \phi_2(a, b; \rho)
\]
Using this property, one can show that:
\[
I_{22} = \rho \cdot I_{11} + \phi_2(a, b; \rho)
\]

\subsubsection{Assembling the Final Result}
Substituting the expression for $I_{22}$ back into the decomposed expectation:
\[
\E[\dots] = \sigma_1\sigma_2 (\rho I_{11} + \phi_2(a,b;\rho)) + \mu_1\sigma_2 I_{12} + \mu_2\sigma_1 I_{21} + \mu_1\mu_2 I_{11}
\]
Finally, grouping the terms by $I_{11}$:
\[
\E[\dots] = (\mu_1\mu_2 + \rho\sigma_1\sigma_2)I_{11} + \mu_1\sigma_2 I_{12} + \mu_2\sigma_1 I_{21} + \sigma_1\sigma_2 \phi_2(a,b;\rho)
\]
This is the complete formula as presented in the main section.

\subsection{Rigorous Calculation of $I_{22}$ using Price's Theorem}
Your intuition is excellent. Price's theorem is precisely the fundamental theoretical tool that allows for an elegant solution to this type of problem.

\paragraph{Application of Price's Theorem}
Price's theorem, in its version adapted for standardized variables with a correlation $\rho$ (which is the covariance $\sigma_{12}$ since $\sigma_1 = \sigma_2 = 1$), is stated as follows:
\[
\frac{\partial \E[g(Z_1, Z_2)]}{\partial \rho} = \E\left[\frac{\partial^2 g(Z_1, Z_2)}{\partial z_1 \partial z_2}\right]
\]
The expectation $\E[\cdot]$ is calculated over the entire bivariate normal distribution $\phi_2(z_1, z_2; \rho)$.

\paragraph{1. Defining the correct function g}
Our goal is to calculate $I_{22} = \int_a^\infty \int_b^\infty z_1 z_2 \cdot \phi_2(z_1, z_2; \rho) \dd z_1 \dd z_2$.
This integral is, by definition, the expectation of the following function:
\[
g(z_1, z_2) = z_1 z_2 \cdot \indic(z_1 > a) \cdot \indic(z_2 > b)
\]
where $\indic(\cdot)$ is the indicator function. Thus, $I_{22} = \E[g(Z_1, Z_2)]$.

\paragraph{2. Calculating the second derivative of g}
Calculating the second derivative $\frac{\partial^2 g}{\partial z_1 \partial z_2}$ would require the theory of distributions (with Dirac delta functions) due to the discontinuities from the indicator functions.

However, we can use a trick by applying the theorem to a simpler function. Consider the derivative of $I_{11}$ with respect to $\rho$, which is a special case of this theorem:
\[
\frac{\partial I_{11}}{\partial \rho} = \frac{\partial \E[\indic(z_1 > a) \cdot \indic(z_2 > b)]}{\partial \rho}
\]
A rigorous calculation (involving differentiation under the integral sign and integration by parts) shows that:
\[
\frac{\partial I_{11}}{\partial \rho} = \phi_2(a, b; \rho)
\]
This is the most famous and directly usable result stemming from this principle.

\paragraph{3. Connecting this result to $I_{22}$}
There is another fundamental identity, which can be proven by a simple integration by parts:
\[
\int_a^\infty z_1 \phi_2(z_1, z_2; \rho) \dd z_1 = \rho z_2 \int_a^\infty \phi_2(z_1, z_2; \rho) \dd z_1 + (1-\rho^2)\phi_2(a, z_2; \rho)
\]
This relationship connects the first-order moment to the zero-order moment. While integrating this identity provides a path, the most direct result from these principles is the remarkably simple formula for $I_{22}$:
\[
I_{22} = \rho \cdot I_{11} + \phi_2(a, b; \rho)
\]
This formula is the "missing link" we need. It expresses the truncated cross-moment ($I_{22}$) in terms of the truncated probability ($I_{11}$) and the value of the probability density function at the truncation point. Your intuition to use Price's theorem was the key to justifying the existence of such a relationship.

\paragraph{Conclusion of the Calculation}
With this formula for $I_{22}$, we can finalize the calculation of the total expectation as shown previously.

\subsection{Detailed Calculation of an Auxiliary Integral}
During the discussion, we performed a detailed calculation for an integral term. Let's retrace the steps. The integral in question was:
\[
I = - \frac{\rho}{\sqrt{1-\rho^2}} \int_b^\infty \phi_1(z_2) \phi_1\left(\frac{a+\rho z_2}{\sqrt{1-\rho^2}}\right) \dd z_2
\]

\paragraph{Step 1: Rewriting the integrand}
We use the property that links the bivariate PDF $\phi_2$ to the product of univariate PDFs:
\[
\phi_2(x, y; \rho) = \phi_1(y) \frac{1}{\sqrt{1-\rho^2}} \phi_1\left(\frac{x - \rho y}{\sqrt{1-\rho^2}}\right)
\]
Our integrand is $\phi_1(z_2) \phi_1\left(\frac{a+\rho z_2}{\sqrt{1-\rho^2}}\right)$. Since $\phi_1$ is an even function, $\phi_1(u) = \phi_1(-u)$. Therefore:
\[
\phi_1\left(\frac{a+\rho z_2}{\sqrt{1-\rho^2}}\right) = \phi_1\left(\frac{-a-\rho z_2}{\sqrt{1-\rho^2}}\right)
\]
This allows us to match the structure of the PDF decomposition. By setting $x=-a$ and $y=z_2$:
\[
\phi_2(-a, z_2; \rho) = \phi_1(z_2) \frac{1}{\sqrt{1-\rho^2}} \phi_1\left(\frac{-a - \rho z_2}{\sqrt{1-\rho^2}}\right)
\]
From this, we deduce the relationship for our integrand:
\[
\phi_1(z_2) \phi_1\left(\frac{a+\rho z_2}{\sqrt{1-\rho^2}}\right) = \sqrt{1-\rho^2} \cdot \phi_2(-a, z_2; \rho)
\]

\paragraph{Step 2: Substitution into the integral}
Replacing this expression in our integral $I$:
\[
I = - \frac{\rho}{\sqrt{1-\rho^2}} \int_b^\infty \left( \sqrt{1-\rho^2} \cdot \phi_2(-a, z_2; \rho) \right) \dd z_2
\]
The $\sqrt{1-\rho^2}$ terms cancel out:
\[
I = -\rho \int_b^\infty \phi_2(-a, z_2; \rho) \dd z_2
\]

\paragraph{Step 3: Calculating the remaining integral}
The remaining integral is $\int_b^\infty \phi_2(-a, z_2; \rho) \dd z_2$. We can again use the conditional density formula, this time conditioning on the first variable:
\[
\phi_2(-a, z_2; \rho) = \phi_1(-a) \frac{1}{\sqrt{1-\rho^2}} \phi_1\left(\frac{z_2 - \rho(-a)}{\sqrt{1-\rho^2}}\right) = \phi_1(a) \frac{1}{\sqrt{1-\rho^2}} \phi_1\left(\frac{z_2 + \rho a}{\sqrt{1-\rho^2}}\right)
\]
The integral becomes:
\[
\int_b^\infty \phi_1(a) \frac{1}{\sqrt{1-\rho^2}} \phi_1\left(\frac{z_2 + \rho a}{\sqrt{1-\rho^2}}\right) \dd z_2 = \frac{\phi_1(a)}{\sqrt{1-\rho^2}} \int_b^\infty \phi_1\left(\frac{z_2 + \rho a}{\sqrt{1-\rho^2}}\right) \dd z_2
\]

\paragraph{Step 4: Change of variable}
Let $u = \frac{z_2 + \rho a}{\sqrt{1-\rho^2}}$. The differential is $\dd u = \frac{\dd z_2}{\sqrt{1-\rho^2}}$, so $\dd z_2 = \sqrt{1-\rho^2}\dd u$. The lower bound becomes $u_{min} = \frac{b + \rho a}{\sqrt{1-\rho^2}}$. The integral transforms to:
\[
\frac{\phi_1(a)}{\sqrt{1-\rho^2}} \int_{\frac{b+\rho a}{\sqrt{1-\rho^2}}}^\infty \phi_1(u) (\sqrt{1-\rho^2} \dd u) = \phi_1(a) \int_{\frac{b+\rho a}{\sqrt{1-\rho^2}}}^\infty \phi_1(u) \dd u
\]

\paragraph{Step 5: Final result}
By definition of the standard normal CDF $\Phi_1$, the integral is $1 - \Phi_1\left(\frac{b+\rho a}{\sqrt{1-\rho^2}}\right) = \Phi_1\left(-\frac{b+\rho a}{\sqrt{1-\rho^2}}\right)$.
So, the result of Step 3 is $\phi_1(a) \Phi_1\left(-\frac{a\rho+b}{\sqrt{1-\rho^2}}\right)$.
Finally, we substitute this back into the expression for $I$ from Step 2:
\[
I = -\rho \left[ \phi_1(a) \Phi_1\left(-\frac{a\rho+b}{\sqrt{1-\rho^2}}\right) \right]
\]

\end{document}